\appendixsection{Theoretical Foundations of HSIC and HIDE}
\label{app:hsic_theory}

\subsection*{A.1 Measuring Statistical Dependence using HSIC}
\label{subsec:hsic_rkhs}
Let's assume that we want to assess the statistical dependence between two random variables $X$ and $Y$, taking values in separable topological spaces $\mathcal{X}$ and $\mathcal{Y}$, respectively. The statistical relationship between $X$ and $Y$ is governed by their joint probability distribution $P_{XY}$. HSIC quantifies this dependence by comparing $P_{XY}$ to the product of the marginals, $P_X P_Y$, within Reproducing Kernel Hilbert Spaces (RKHSs). An RKHS $\mathcal{F}_k$, associated with a kernel $k: \mathcal{X} \times \mathcal{X} \to \mathbb{R}$ allows implicit mapping to a feature space via $\phi: \mathcal{X} \to \mathcal{F}_k$, where $k(x, x') = \langle \phi(x), \phi(x') \rangle_{\mathcal{F}_k}$, with $\langle \cdot, \cdot \rangle_{\mathcal{F}_k}$ denoting the inner product in the RKHS $\mathcal{F}_k$. This mapping facilitates the detection of non-linear dependencies.


% Given two RKHSs, $\mathcal{F}$ with kernel $k_X$ for variable $X$, and $\mathcal{G}$ with kernel $k_Y$ for variable $Y$ (corresponding to feature maps $\phi$ and $\psi$ respectively), the cross-covariance operator $C_{XY}: \mathcal{G} \to \mathcal{F}$ is defined as:
% \begin{equation*}
% \begin{small}
% C_{XY} \coloneqq \mathbb{E}_{XY}[\phi(X) \otimes \psi(Y)] - \mathbb{E}_X[\phi(X)] \otimes \mathbb{E}_Y[\psi(Y)]
% \end{small}
% \end{equation*}
% where $\otimes$ denotes the tensor product. Intuitively, $C_{XY}$ captures the covariance between the representations of $X$ and $Y$ in their respective RKHSs. If $X$ and $Y$ are independent, $C_{XY}$ is the zero operator.
\changedtwo{\revisionr3{Intuitively, standard linear metrics (like cosine similarity) only capture direct geometric alignment. In contrast, evaluating dependence in an RKHS allows us to detect arbitrary, non-linear relationships. The cross-covariance operator, $C_{XY}$, acts as a high-dimensional analog to the standard covariance matrix, capturing the covariance between the transformed representations of $X$ and $Y$ in their respective feature spaces. A reduced cross-covariance magnitude motivates the decoupling hypothesis, but the observed score and proxy correlations do not establish this as a general consequence or cause of hallucination. Formally, given two RKHSs, $\mathcal{F}$ with kernel $k_X$ for variable $X$, and $\mathcal{G}$ with kernel $k_Y$ for variable $Y$ (corresponding to feature maps $\phi$ and $\psi$ respectively), the cross-covariance operator $C_{XY}: \mathcal{G} \to \mathcal{F}$ is defined as:}}
\begin{equation*}
C_{XY} \coloneqq \mathbb{E}_{XY}[\phi(X) \otimes \psi(Y)] - \mathbb{E}_X[\phi(X)] \otimes \mathbb{E}_Y[\psi(Y)]
\end{equation*}
where $\otimes$ denotes the tensor product. If $X$ and $Y$ are independent, $C_{XY}$ is the zero operator.

The Hilbert-Schmidt norm of an operator $C: \mathcal{G} \to \mathcal{F}$, denoted by $\|C\|_{HS}$, is a generalization of the Frobenius norm for matrices and can be defined as $\|C\|_{HS}^2 = \sum_{i,j} |\langle C g_j, f_i \rangle_{\mathcal{F}}|^2$, where $\{f_i\}$ and $\{g_j\}$ are orthonormal bases of $\mathcal{F}$ and $\mathcal{G}$, respectively. It essentially measures the \textit{magnitude} of the operator.


\begin{definition}
\label{def:hsic}
\textit{HSIC is the squared Hilbert-Schmidt norm of the cross-covariance operator $C_{XY}$:}
\begin{equation}
\text{\textit{HSIC}}(X,Y; k_X, k_Y) \coloneqq \|C_{XY}\|^2_{HS}
\end{equation}
\end{definition}

\noindent A crucial property of HSIC is its relationship with statistical independence, formalized as follows:

\begin{lemma}[\citet{gretton05a}]
\label{lemma:hsic_independence_merged}
\textit{Let $k_X$ and $k_Y$ be characteristic kernels on $\mathcal{X}$ and $\mathcal{Y}$, respectively. Then $\operatorname{HSIC}(X,Y; k_X, k_Y) = 0$ if and only if $X$ and $Y$ are statistically independent.}
\end{lemma}
A kernel $k$ is termed characteristic if the mapping from probability measures on its domain to mean elements in its RKHS, $\mu \mapsto \mathbb{E}_{X \sim \mu}[\phi(X)]$, is injective \citep{JMLR:v11:sriperumbudur10a}. This ensures that the mean embedding $\mathbb{E}[\phi(X)]$ uniquely determines $P_X$. For our method, we use the Radial Basis Function (RBF) kernel, $k_{\text{RBF}}(x, x') = \exp(-\gamma \|x - x'\|_2^2)$ for a bandwidth parameter $\gamma > 0$.

\begin{lemma}[\citet{5122}]
\label{lemma:rbf_characteristic_merged}
\textit{The RBF kernel is characteristic on $\mathbb{R}^d$.}
\end{lemma}
The conjunction of Lemma~\ref{lemma:hsic_independence_merged} and Lemma~\ref{lemma:rbf_characteristic_merged} establishes HSIC, when computed with RBF kernels, as a robust and theoretically sound metric for detecting statistical independence, forming the cornerstone of our approach.

\subsection*{A.2 Empirical Estimation of HSIC}
\label{subsec:empirical_hsic_foundation}
In practice, we rely on empirical estimates of HSIC to quantify statistical independence from a finite set of $n$ samples. Given samples $\{x_1, \dots, x_n\}$ from the distribution of $X$ and $\{y_1, \dots, y_n\}$ from the distribution of $Y$, we form $n \times n$ Gram matrices $\mathbf{K}_X$ and $\mathbf{K}_Y$, where $\mathbf{K}_{X,ij} = k_X(x_i, x_j)$ and $\mathbf{K}_{Y,ij} = k_Y(y_i, y_j)$.

A common, though biased, empirical estimator for HSIC is \citep{gretton05a}:

\begin{equation}
\widehat{\operatorname{HSIC}}_b(X,Y) = \frac{1}{(n-1)^2} \text{Tr}(\mathbf{K}_X \mathbf{H} \mathbf{K}_Y \mathbf{H})
\label{eq:biased_HSIC_foundation}
\end{equation}
where $\mathbf{H} = \mathbf{I} - \frac{1}{n}\mathbf{1}\mathbf{1}^T$ is the $n \times n$ centering matrix, $\mathbf{I}$ is the identity matrix, and $\mathbf{1}$ is an $n \times 1$ column vector of ones.

For improved performance, particularly with smaller $n$, an unbiased estimator is preferred \citep{JMLR:v13:song12a}. \revisionr3{The unbiased estimator is given by:}
\begin{multline}
\widehat{\operatorname{HSIC}}_u(X,Y) = \frac{1}{n(n-3)} \bigg[ \text{Tr}(\tilde{\mathbf{K}}_X\tilde{\mathbf{K}}_Y) +   \frac{(\mathbf{1}^T \tilde{\mathbf{K}}_X \mathbf{1}) (\mathbf{1}^T \tilde{\mathbf{K}}_Y \mathbf{1})}{(n-1)(n-2)} -  \frac{2}{n-2} \mathbf{1}^T \tilde{\mathbf{K}}_X\tilde{\mathbf{K}}_Y \mathbf{1} \bigg]
\label{eq:unbiased_HSIC_foundation}
\end{multline}
Here, $n$ is the number of samples. The matrices $\tilde{\mathbf{K}}_X$ and $\tilde{\mathbf{K}}_Y$ are derived from $\mathbf{K}_X$ and $\mathbf{K}_Y$, respectively by setting their diagonal entries to zero: $\tilde{\mathbf{K}}_{X,ij} = (1-\delta_{ij})\mathbf{K}_{X,ij}$, where $\delta_{ij}$ is the Kronecker delta  (1 if $i=j$, and 0 otherwise). This estimator requires $n \ge 4$ for its standard definition.

\subsection*{A.3 Properties of the \method~Score}
\label{sec:properties-hide}
We now analyze the properties of our proposed \method~score (Equation ~\ref{eq:hide-score}), denoted as \(\widehat{\operatorname{HSIC}}_{\method}\), \revisionr3{which adapts the standard unbiased HSIC estimator} to be computable for any number of effective tokens \(n \equiv n_{\text{eff}} \ge 1\).

Let \(\{(x_i,y_i)\}_{i=1}^{n}\) be the set of hidden state representations for the selected input and output tokens. We use the RBF kernel, which is bounded such that \(0 < k(u,v) \le 1\) for any vectors \(u,v\). The score is computed using Gram matrices with zeroed diagonals, \(\tilde{\mathbf{K}}_X\) and \(\tilde{\mathbf{K}}_Y\), as defined in Equation~\ref{eq:unbiased_HSIC_foundation}.

\begin{lemma}[\textbf{Behaviour for Small Sample Sizes}]
\label{lem:small_n_properties}
\textit{The \(\widehat{\operatorname{HSIC}}_{\method}\) estimator is computable for small sample sizes where \(\widehat{\operatorname{HSIC}}_u\) is undefined. Specifically:}
\begin{enumerate}
    \item[\textit{(i)}] \textit{If \(n=1\), then \(\widehat{\operatorname{HSIC}}_{\method}(X,Y) = 0\).}
    \item[\textit{(ii)}] \textit{If \(n=2\), then \(\widehat{\operatorname{HSIC}}_{\method}(X,Y) = \frac{1}{4} k_X(x_1,x_2)k_Y(y_1,y_2)\), which implies \(0 < \widehat{\operatorname{HSIC}}_{\method}(X,Y) \le 1/4\).}
\end{enumerate}
\end{lemma}

\begin{proof}
The proof of the above lemma follows from direct substitution into Equation~\ref{eq:hide-score} for computing \(\widehat{\operatorname{HSIC}}_{\method}\).

(i) For \(n=1\), the zero-diagonal matrices are \(\tilde K_X = [0]\) and \(\tilde K_Y = [0]\). All terms in Equation~\ref{eq:hide-score} thus becomes zero, implying \(\widehat{\operatorname{HSIC}}_{\method} = 0\).

(ii) For \(n=2\), let \(a = k_X(x_1,x_2)\) and \(b = k_Y(y_1,y_2)\). The matrices are, therefore, \(\tilde K_X = \begin{pmatrix} 0 & a \\ a & 0 \end{pmatrix}\) and \(\tilde K_Y = \begin{pmatrix} 0 & b \\ b & 0 \end{pmatrix}\). The constituent terms of Equation~\ref{eq:hide-score} are: \(\text{Tr}(\tilde K_X\tilde K_Y) = 2ab\); \(\mathbf{1}^{\!\top}\tilde K_X\mathbf{1} = 2a\); \(\mathbf{1}^{\!\top}\tilde K_Y\mathbf{1} = 2b\); and \(\mathbf{1}^{\!\top}\tilde K_X\tilde K_Y\mathbf{1} = 2ab\). Substituting these into Equation~\ref{eq:hide-score} for \(n=2\) yields:
\[ \widehat{\operatorname{HSIC}}_{\method}(X,Y) = \frac{1}{2^2} \left[ 2ab + \frac{(2a)(2b)}{2^2} - \frac{2}{2}(2ab) \right] = \frac{1}{4} [2ab + ab - 2ab] = \frac{ab}{4}. \]
Since the RBF kernel ensures \(a, b \in (0,1]\) for distinct points, the result \( (0, 1/4]\) follows. This confirms the score yields a deterministic, non-negative value for \(n=2\).
\end{proof}


\begin{lemma}[\textbf{Asymptotic Properties}]
\label{lem:asymptotic_properties}
\revisionr3{\textit{For independent and identically distributed paired observations $(x_i,y_i)$ and fixed bounded RBF kernels, the \(\widehat{\operatorname{HSIC}}_{\method}\) estimator is asymptotically unbiased and strongly consistent.}}
\begin{enumerate}
    \item[\textit{(i)}] \textit{(Asymptotic Bias) The bias of the estimator vanishes as \(n \to \infty\):
    \[ \bigl|\mathbb{E}[\widehat{\operatorname{HSIC}}_{\method}] - \operatorname{HSIC}(X,Y)\bigr| = O(1/n). \]}
    \item[\textit{(ii)}] \textit{(Strong Consistency) As \(n \to \infty\), \(\widehat{\operatorname{HSIC}}_{\method}(X,Y)\) converges almost surely to the true population HSIC:
    \[ \widehat{\operatorname{HSIC}}_{\method}(X,Y) \xrightarrow{\text{a.s.}} \operatorname{HSIC}(X,Y). \]}
\end{enumerate}
\end{lemma}

% \appendixsection{Proofs of Lemmas}
% \label{appendix:proofs}
% We discuss detailed mathematical proofs of Lemmas 3 and 4, introduced in \cref{sec:properties-hide}, below:
% \\
% \noindent\textbf{Lemma 3 (Behaviour for Small Sample Sizes).}
% \label{lem:small_n_properties}
% \textit{The \(\widehat{\operatorname{HSIC}}_{\method}\) estimator is computable for small sample sizes where \(\widehat{\operatorname{HSIC}}_u\) is undefined. Specifically:}
% \begin{enumerate}
%     \item[\textit{(i)}] \textit{If \(n=1\), then \(\widehat{\operatorname{HSIC}}_{\method}(X,Y) = 0\).}
%     \item[\textit{(ii)}] \textit{If \(n=2\), then \(\widehat{\operatorname{HSIC}}_{\method}(X,Y) = \frac{1}{4} k_X(x_1,x_2)k_Y(y_1,y_2)\), which implies \(0 < \widehat{\operatorname{HSIC}}_{\method}(X,Y) \le 1/4\).}
% \end{enumerate}

% \begin{proof}
% The proof of the above lemma follows from direct substitution into Equation~\ref{eq:hide-score} for computing \(\widehat{\operatorname{HSIC}}_{\method}\).

% (i) For \(n=1\), the zero-diagonal matrices are \(\tilde K_X = [0]\) and \(\tilde K_Y = [0]\). All terms in Equation~\ref{eq:hide-score} thus becomes zero, implying \(\widehat{\operatorname{HSIC}}_{\method} = 0\).

% (ii) For \(n=2\), let \(a = k_X(x_1,x_2)\) and \(b = k_Y(y_1,y_2)\). The matrices are, therefore, \(\tilde K_X = \begin{pmatrix} 0 & a \\ a & 0 \end{pmatrix}\) and \(\tilde K_Y = \begin{pmatrix} 0 & b \\ b & 0 \end{pmatrix}\). The constituent terms of Equation~\ref{eq:hide-score} are: \(\text{Tr}(\tilde K_X\tilde K_Y) = 2ab\); \(\mathbf{1}^{\!\top}\tilde K_X\mathbf{1} = 2a\); \(\mathbf{1}^{\!\top}\tilde K_Y\mathbf{1} = 2b\); and \(\mathbf{1}^{\!\top}\tilde K_X\tilde K_Y\mathbf{1} = 2ab\). Substituting these into Equation~\ref{eq:hide-score} for \(n=2\) yields:
% \[ \widehat{\operatorname{HSIC}}_{\method}(X,Y) = \frac{1}{2^2} \left[ 2ab + \frac{(2a)(2b)}{2^2} - \frac{2}{2}(2ab) \right] = \frac{1}{4} [2ab + ab - 2ab] = \frac{ab}{4}. \]
% Since the RBF kernel ensures \(a, b \in (0,1]\) for distinct points, the result \( (0, 1/4]\) follows. This confirms the score yields a deterministic, non-negative value for \(n=2\).
% \end{proof}

% \noindent\textbf{Lemma 4 (Asymptotic Properties).}
% \label{lem:asymptotic_properties}
% \revisionr3{\textit{For independent and identically distributed paired observations $(x_i,y_i)$ and fixed bounded RBF kernels, the \(\widehat{\operatorname{HSIC}}_{\method}\) estimator is asymptotically unbiased and strongly consistent.}}
% \begin{enumerate}
%     \item[\textit{(i)}] \textit{(Asymptotic Bias) The bias of the estimator vanishes as \(n \to \infty\):
%     \[ \bigl|\mathbb{E}[\widehat{\operatorname{HSIC}}_{\method}] - \operatorname{HSIC}(X,Y)\bigr| = O(1/n). \]}
%     \item[\textit{(ii)}] \textit{(Strong Consistency) As \(n \to \infty\), \(\widehat{\operatorname{HSIC}}_{\method}(X,Y)\) converges almost surely to the true population HSIC:
%     \[ \widehat{\operatorname{HSIC}}_{\method}(X,Y) \xrightarrow{\text{a.s.}} \operatorname{HSIC}(X,Y). \]}
% \end{enumerate}

\begin{proof}
\revisionr3{Write $A_n=\operatorname{Tr}(\tilde K_X\tilde K_Y)$,
$B_n=(\mathbf1^\top\tilde K_X\mathbf1)(\mathbf1^\top\tilde K_Y\mathbf1)$,
and $C_n=\mathbf1^\top\tilde K_X\tilde K_Y\mathbf1$.
Since the kernel entries are bounded by one, $|A_n|\leq n^2$,
$|B_n|\leq n^4$, and $|C_n|\leq n^3$.
The differences between the coefficients of these three terms in the adapted
and unbiased estimators satisfy, as $n\to\infty$,}
\revisionr3{\begin{align*}
\left|\frac{1}{n^2}-\frac{1}{n(n-3)}\right|&=O(n^{-3}),\\
\left|\frac{1}{n^4}-\frac{1}{n(n-3)(n-1)(n-2)}\right|&=O(n^{-5}),\\
\left|\frac{2}{n^3}-\frac{2}{n(n-3)(n-2)}\right|&=O(n^{-4}).
\end{align*}}
\revisionr3{Multiplying these bounds gives
$|\widehat{\operatorname{HSIC}}_{\method}-\widehat{\operatorname{HSIC}}_u|=O(1/n)$,
uniformly over the observations. Since the standard estimator is unbiased
under the stated sampling assumptions, taking expectations proves (i).
Its strong consistency under the same assumptions, together with the vanishing
uniform difference, proves (ii). This limiting result does not imply unbiasedness
at a fixed small token count, nor does it establish independence for token
representations selected from a single generated answer.}
\end{proof}
